%-*-tex-*-
% Copyright Michael J. Ferguson, INRS-Telecommunications
% All rights reserved. 

% ====== Multiple Language Support =========
% This is done via an internal change to TeX itself or through the
% introduction of a counter called \language. If Bilingual TeX exists,
% then \language has been defined. 

\ifundefined{language} \n@ewcount\language \language=0 \fi 

\def\englishversion{\input english \relax}
\def\versionfrancaise{\input francais \relax}



%======= are some commands for << and >> in French ======
% ======== can be made neater when we modify fonts
% \newcount\language % ..... this is necessary if NOT running Multilingual TeX
\language = 0   % 0-english 1-french ...
\def\d@qf{\raise .5ex\hbox{$\scriptscriptstyle{\langle\!\langle}$}\kern.25em}
\def\g@qf{\kern.25em\raise .5ex\hbox{$\scriptscriptstyle{\rangle\!\rangle}$}}
\def\ldq{\ifcase\language
                \lq\lq \else\d@qf\fi} %changes < into `` in english
\def\rdq{\ifcase\language
                          \rq\rq\else\g@qf\fi} %changes > into '' in english